\documentclass[11pt]{article}

\usepackage[utf8]{inputenc}
\usepackage[T1]{fontenc}
\usepackage{amsmath, amssymb, amsthm}
\usepackage{geometry}
\geometry{margin=1in}
\usepackage{hyperref}
\hypersetup{
  colorlinks=true,
  linkcolor=blue,
  citecolor=blue,
  urlcolor=blue,
}
\usepackage{mathtools}

\newtheorem{definition}{Definition}[section]
\newtheorem{theorem}[definition]{Theorem}
\newtheorem{lemma}[definition]{Lemma}
\newtheorem{proposition}[definition]{Proposition}
\newtheorem{conjecture}[definition]{Conjecture}
\newtheorem{remark}[definition]{Remark}
\newtheorem{example}[definition]{Example}

\newcommand{\Rnnext}{\overline{\mathbb{R}}_{\ge 0}}
\newcommand{\Rnn}{\mathbb{R}_{\ge 0}}
\newcommand{\Pcal}{\mathcal{P}}
\newcommand{\Acal}{\mathcal{A}}
\newcommand{\Scal}{\mathcal{S}}
\newcommand{\Omk}{\Omega_k}
\newcommand{\Omsf}{\Omega_k^{\mathrm{sf}}}
\newcommand{\Omlf}{\Omega_k^{\mathrm{lf}}}
\newcommand{\Omcap}{\Omega_k^{\cap}}

\title{The Sheaf-Theoretic Framework for the Constructive\\
  Yang--Mills Mass Gap Program:\\
  Sites, Bound-Witnesses, and Summability Certificates}
\author{Arthur Petron}
\date{\today}

\begin{document}
\maketitle

\begin{abstract}
We define the sheaf-theoretic framework underlying the constructive Yang--Mills mass gap program: the configuration site over $\Omk = G^{E_k}$ with its three-element small-field/large-field/buffer cover, the bound-witness presheaf $\Acal_k$ valued in $(\Rnnext, \le)$, its flabbiness, and the summability-certificate presheaf $\Scal_k$ carrying Kotecký--Preiss data $(\rho, K)$.  We prove that $\Acal_k$ is flabby in the domination sense; the proof is constructive (no excluded middle), exhibiting the constant-$\infty$ section as the trivial dominator.  Flabbiness is what makes the sheaf condition automatic for $\Acal_k$, isolating the analytic content of the framework in the compatibility condition between $\Acal_k$ and $\Scal_k$, which we state as a named definition.  This document is the formal reference companion to \texttt{docs/sheaf-framework.md} (prose argument) and \texttt{lean/YMGap/Sheaf/} (Lean formalization).  It is group-independent: no $SU(2)$-specific constants appear.  Its scope corresponds to Milestone~2.1 of \texttt{work\_plan.md}.
\end{abstract}

\tableofcontents

\section{Introduction}
\label{sec:intro}

The constructive Yang--Mills mass gap program (\texttt{docs/outline.md}) reorganizes the small-field/large-field obstruction of the classical Bałaban-style construction as a \emph{cocycle problem rather than a restriction problem}: the small-field and large-field activity bounds are not restrictions of a single function to two subsets of $\Omk$, but two distinct upper bounds on the same activity derived under two different approximations of the action (\texttt{docs/context.md}, \S3.1).  The gluing of these bounds on their overlap is therefore by \emph{domination} rather than \emph{equality}, and the question of whether the two bounds are simultaneously usable becomes a question of a coupled pair of presheaves on a shared site.

This document fixes the categorical and analytic vocabulary of that framework, in the scope required by Milestone~2.1 of \texttt{work\_plan.md}.  We do not in this document close any analytic obligation; the constants for the three regimes are computed in Milestones 3.1--3.3 against the framework defined here.

\section{The configuration site}
\label{sec:site}

Throughout, $G$ is a compact simple Lie group, $k \in \mathbb{N}$ an RG scale, $a_k > 0$ the lattice spacing at scale $k$, $E_k$ the set of (oriented) lattice edges at scale $k$, and $\Omk := G^{E_k}$ the link-configuration space.  We do not specialize $G$: this section is group-independent.

\begin{definition}[Polymer]
\label{def:polymer}
A \emph{polymer} is a finite, connected subgraph $X$ of the plaquette adjacency graph at scale $k$.  We write $\Pcal_k$ for the set of polymers and $|X|$ for the plaquette count of $X \in \Pcal_k$.  Polymer inclusion $X \subseteq X'$ is the inclusion of underlying plaquette sets.
\end{definition}

\begin{definition}[Three-element configuration cover]
\label{def:three-cover}
Let $\delta_k > 0$ be the smallness parameter (computable from $g_k$ via $\delta_k \sim g_k^\alpha$ for some $\alpha \in (0,1)$; see outline \S6.4).  The three-element cover of $\Omk$ is
\begin{align*}
  \Omsf  &:= \{ U \in \Omk : \forall p,\ d(U_p, 1) < \delta_k \},\\
  \Omlf  &:= \{ U \in \Omk : \exists p,\ d(U_p, 1) > \delta_k / 2 \},\\
  \Omcap &:= \{ U \in \Omk : \exists p,\ \delta_k / 2 < d(U_p, 1) < 2 \delta_k \},
\end{align*}
where $d(\cdot, 1)$ is the geodesic distance on $G$ from the identity.  By construction $\Omsf \cup \Omlf \cup \Omcap = \Omk$ and the buffer $\Omcap$ contains the pairwise overlap $\Omsf \cap \Omlf$ but is strictly larger: $\Omcap$ absorbs every plaquette in the indeterminate annulus $\delta_k / 2 < d(U_p, 1) < 2 \delta_k$, not only the strict intersection.

\textbf{Critical remark.}  This is a \emph{three-element} cover.  The buffer $\Omcap$ is a first-class open subset, not a derived "overlap" between two regions.  The cocycle observation of \texttt{docs/context.md} \S3.1 forces this: two distinct local bounds (small-field and large-field) need to be \emph{reconciled} on $\Omcap$, not glued by equality, so $\Omcap$ must carry sections of its own.
\end{definition}

\begin{definition}[Configuration site $\mathrm{Site}(\Omk)$]
\label{def:site}
The \emph{configuration site at scale $k$} is the Grothendieck site whose
\begin{itemize}
  \item \emph{objects} are pairs $(U, X)$ with $U \subseteq \Omk$ an open subset (in any topology on $\Omk$ admitting decidable membership for the cover of Definition~\ref{def:three-cover}) and $X \in \Pcal_k$ a polymer;
  \item \emph{morphisms} $(U, X) \to (U', X')$ are pairs $(\iota, \jmath)$ where $\iota \colon U \hookrightarrow U'$ is the inclusion of opens and $\jmath \colon X \subseteq X'$ is the inclusion of polymers (so $U \subseteq U'$ and $X \subseteq X'$ as sets);
  \item \emph{coverage} is generated by the single covering family
  \begin{equation}
    \label{eq:cover}
    \big\{ (\Omsf, X),\ (\Omlf, X),\ (\Omcap, X) \big\} \;\Rightarrow\; (\Omk, X), \qquad X \in \Pcal_k.
  \end{equation}
\end{itemize}
\end{definition}

A presheaf on $\mathrm{Site}(\Omk)$ is a functor $\mathrm{Site}(\Omk)^{\mathrm{op}} \to \mathbf{Set}$; for our purposes we only need presheaves valued in partially ordered sets, with restriction maps that preserve the order.

\section{The bound-witness presheaf $\Acal_k$}
\label{sec:awitness}

We use the extended nonneg reals $\Rnnext := [0, \infty]$ as the value type, ordered by the standard $\le$.  We write $\sup$ for binary sup and treat $+\infty$ as the order-theoretic top.

\paragraph{Constructive reading.}  Where below we write "$f \colon U \to \Rnnext$", the intended reading throughout is Bishop-constructive: each value $f(u)$ is presented together with a computable modulus of convergence (a Cauchy modulus when finite, or an explicit witness $f(u) = \infty$).  This is the reading recommended by Milestone~2.1's open-question paragraph in \texttt{work\_plan.md}, and the level at which downstream Lean instantiations should be made.  The structural lemmas of this section do not depend on the constructive interpretation, but their downstream use does.

\begin{definition}[Polymer activity]
\label{def:activity}
We assume given a polymer activity $|a_k(X)| \colon \Omk \to \Rnnext$ for each $X \in \Pcal_k$, measurable and bounded.  Specific forms (small-field perturbative activity, large-field heat-kernel activity, and their reconciliation on the buffer) are pinned in Milestones 1.3, 1.4, and 3.1--3.3.  Here $|a_k(X)|$ is an abstract parameter.
\end{definition}

\begin{definition}[Bound-witness presheaf $\Acal_k$]
\label{def:Akpresheaf}
For each object $(U, X)$ of the configuration site, set
\begin{equation*}
  \Acal_k(U, X) \;:=\; \big\{ f \colon U \to \Rnnext \;\big|\; |a_k(X)|(u) \le f(u) \ \text{for every}\ u \in U \big\}.
\end{equation*}
For a morphism $(U, X) \to (U', X')$ (so $U \subseteq U'$, $X \subseteq X'$), define the \emph{restriction map}
\begin{equation*}
  \rho_{(U', X') \to (U, X)} \colon \Acal_k(U', X') \to \Acal_k(U, X), \qquad f \mapsto f|_U.
\end{equation*}
The set $\Acal_k(U, X)$ is partially ordered by pointwise inequality on $U$: $f \le g$ iff $f(u) \le g(u)$ for every $u \in U$.  In this order, larger means \emph{looser bound}.
\end{definition}

\begin{lemma}[Restriction stability]
\label{lem:restriction-stability}
The map $\rho_{(U', X') \to (U, X)}$ is well-defined: if $f \in \Acal_k(U', X')$ then $f|_U \in \Acal_k(U, X)$.
\end{lemma}

\begin{proof}
For $u \in U \subseteq U'$ we have $|a_k(X')|(u) \le f(u)$ by the admissibility of $f$ on $U'$.  Since $X \subseteq X'$ as polymers, the activity $|a_k(X)|$ is dominated by $|a_k(X')|$ on $u$ (the larger polymer's activity dominates the smaller's; see Definition~\ref{def:activity}'s standing assumption).  Hence $|a_k(X)|(u) \le |a_k(X')|(u) \le f(u)$, and $f|_U$ is admissible.
\end{proof}

\begin{remark}[The order is the analytic data]
The set $\Acal_k(U, X)$ is an \emph{upper set} in $\Rnnext^U$: if $f$ is admissible and $g \ge f$ pointwise, then $g$ is also admissible.  In particular, $\Acal_k(U, X)$ is nonempty as soon as some bound exists --- and the constant function $\infty$ on $U$ is always in $\Acal_k(U, X)$.
\end{remark}

\subsection{Flabbiness}
\label{sec:flabby}

\begin{theorem}[Flabbiness of $\Acal_k$]
\label{thm:flabbiness}
The bound-witness presheaf $\Acal_k$ is \emph{flabby in the domination sense}: for every morphism $(U, X) \to (V, X)$ in the configuration site and every $s \in \Acal_k(U, X)$, there exists $t \in \Acal_k(V, X)$ with $s \le t|_U$ in the pointwise order on $\Rnnext^U$.
\end{theorem}

\begin{proof}
Define $t \colon V \to \Rnnext$ by $t \equiv \infty$ (constant function).  Then for every $v \in V$, $|a_k(X)|(v) \le \infty = t(v)$, so $t \in \Acal_k(V, X)$.  Restricting to $U$: $t|_U(u) = \infty \ge s(u)$ for every $u \in U$, so $s \le t|_U$.
\end{proof}

\begin{remark}[Constructive content of the proof]
\label{rem:flabby-constructive}
The proof of Theorem~\ref{thm:flabbiness} is constructive in the strong sense.  It does not case-split on whether $u$ lies in $U$ or $V \setminus U$ (so it does not need decidability of open membership), it does not invoke excluded middle for the inequality $|a_k(X)|(v) \le \infty$ (this is the structural fact that $\infty$ is the top of $\Rnnext$, available constructively in any presentation), and the witness $t$ is exhibited explicitly as the constant section.

This is by design: the Lean formalization in \texttt{lean/YMGap/Sheaf/BoundWitness.lean} mirrors the same proof, with $\Rnnext$ replaced by an abstract type $T$ carrying a partial order, a top element $\top$, and a binary sup.  Concrete value-type instantiations (rational surrogates; constructive reals in the sense of \texttt{lean/README.md}, when that module lands; the extended nonneg reals here) all enjoy the same flabbiness proof.
\end{remark}

\begin{remark}[Why "domination" rather than "equality"]
\label{rem:domination-flabby}
The standard classical statement of flabbiness, "the restriction map $\Acal_k(V, X) \to \Acal_k(U, X)$ is surjective onto $s$", would require constructing a $t \in \Acal_k(V, X)$ with $t|_U = s$ exactly.  In the framework of \texttt{docs/context.md} \S3.2 this strict version is also true (by piecewise extension: $t|_U := s$, $t|_{V \setminus U} := \infty$), but only \emph{given} decidable membership in $U$.  The domination form chosen here is what the program actually uses downstream: it suffices for the structural content of the sheaf condition --- "every local bound has a global bound dominating it" --- without the decidability hypothesis.
\end{remark}

\begin{remark}[What flabbiness buys us]
\label{rem:flabby-buys}
Flabbiness in the domination sense means the only structural constraint on $\Acal_k$ is the existence of \emph{some} global bound, which is trivially $t \equiv \infty$.  The sheaf condition on $\Acal_k$ is therefore automatic and carries no analytic content.  The analytic content of the framework lives in the \emph{summability constraint} on $\Acal_k$, which we formalize as a second presheaf $\Scal_k$ and a named compatibility condition.
\end{remark}

\section{The summability-certificate presheaf $\Scal_k$}
\label{sec:scertif}

The polymer expansion of the partition function at scale $k$,
\begin{equation*}
  Z_k = \sum_{\{X_1, \ldots, X_n\}} \prod_i a_k(X_i)
\end{equation*}
(sum over compatible finite polymer families), converges absolutely under a Kotecký--Preiss (KP) summability condition: there exist $\rho, K > 0$ such that
\begin{equation}
  \label{eq:kp}
  \sum_{X \in \Pcal_k,\ X \ni x} |a_k(X)|(u)\, e^{\rho |X|} \;\le\; K
  \qquad \text{for all } u \in U,\ x \in X(u),
\end{equation}
where $X \ni x$ denotes that the polymer $X$ contains the plaquette or vertex $x$.  The data $(\rho, K)$ is what we record in $\Scal_k$.

\begin{definition}[Summability-certificate presheaf $\Scal_k$]
\label{def:Skpresheaf}
For each object $(U, X)$ of the configuration site, set
\begin{equation*}
  \Scal_k(U, X) \;:=\; \big\{ (\rho, K) \in \mathbb{R}_{> 0} \times \mathbb{R}_{\ge 0} \;\big|\; (\rho, K) \text{ is a KP datum at } (U, X) \big\}
\end{equation*}
where "KP datum at $(U, X)$" means equation~\eqref{eq:kp} holds for the polymer activity restricted to $U$ and summed over polymers $X' \in \Pcal_k$ with $X' \supseteq X$ (i.e. the polymers extending the marker $X$).  Restriction along a morphism $(U, X) \to (U', X')$ takes $(\rho, K)$ to itself --- the same datum certifies summability on the smaller open $U \subseteq U'$ over the smaller family of polymers $\{X'' : X'' \supseteq X'\} \subseteq \{X'' : X'' \supseteq X\}$.
\end{definition}

The three regimes carry their own KP witnesses, one per element of the configuration cover (Definition~\ref{def:three-cover}):
\begin{equation*}
  (\rho_{\mathrm{sf}}, K_{\mathrm{sf}}) \in \Scal_k(\Omsf, X), \quad
  (\rho_{\mathrm{lf}}, K_{\mathrm{lf}}) \in \Scal_k(\Omlf, X), \quad
  (\rho_{\cap}, K_{\cap}) \in \Scal_k(\Omcap, X).
\end{equation*}
Their specific values for $SU(2)$ are computed in Milestones 3.1--3.3.

\begin{remark}[Why $\Scal_k$ is not automatically nonempty]
Unlike $\Acal_k$, the presheaf $\Scal_k$ has no trivial section: the constant $K \equiv \infty$ does not satisfy equation~\eqref{eq:kp} for any meaningful $\rho > 0$, and finite $K$'s require the activity to actually decay exponentially.  Existence of $\Scal_k$ sections is therefore the analytic content of the framework.  This is the structural opposite of $\Acal_k$'s flabbiness (Remark~\ref{rem:flabby-buys}): $\Acal_k$ is structurally trivial and analytically rich, $\Scal_k$ is structurally simple but analytically demanding.
\end{remark}

\section{Compatibility: where the analytic content lives}
\label{sec:compat}

\begin{definition}[Compatibility between $\Acal_k$ and $\Scal_k$]
\label{def:compatibility}
Let $(U, X)$ be an object of $\mathrm{Site}(\Omk)$.  A bound $s \in \Acal_k(U, X)$ and a KP datum $(\rho, K) \in \Scal_k(U, X)$ are \emph{compatible}, written
\begin{equation*}
  s \rightleftharpoons (\rho, K),
\end{equation*}
if they jointly certify polymer summability at $(U, X)$: for every $u \in U$ and every plaquette $x \in U_p$,
\begin{equation}
  \label{eq:compatibility}
  \sum_{X' \in \Pcal_k,\ X' \supseteq X,\ X' \ni x} s_{X'}(u)\, e^{\rho |X'|} \;\le\; K,
\end{equation}
where $s_{X'} \in \Acal_k(U, X')$ ranges over a coherent family of bound-witness sections (one per polymer $X'$ extending $X$).  Coherence here means each $s_{X'}$ is admissible on $(U, X')$ in the sense of Definition~\ref{def:Akpresheaf}.
\end{definition}

This is the named definition Milestone~2.1 requires (work\_plan.md L191--195).  It is the place where the analytic content of the framework lives: compatibility binds an admissible activity bound (from $\Acal_k$) to a Kotecký--Preiss decay/amplitude pair (from $\Scal_k$), and the inductive step of the program (outline \S10, work\_plan.md Phase~3) is exactly the assertion that compatible families persist along the block-spin pushforward up to a controlled degradation.

\begin{remark}[Where compatibility is closed]
\label{rem:where-closed}
The compatibility predicate~\eqref{eq:compatibility} is not closed in this document for any concrete regime.  Regime-by-regime, it is closed in:
\begin{itemize}
  \item Milestone~3.1 for the small-field regime, $(s^{\mathrm{sf}}, (\rho_{\mathrm{sf}}, K_{\mathrm{sf}}))$ on $(\Omsf, X)$;
  \item Milestone~3.2 for the large-field regime, $(s^{\mathrm{lf}}, (\rho_{\mathrm{lf}}, K_{\mathrm{lf}}))$ on $(\Omlf, X)$;
  \item Milestone~3.3 for the overlap, by combining 3.1 and 3.2 with the framework-level Proposition~\ref{prop:summability-preservation} of \S\ref{sec:buffer-summability}.  The framework-level claim (that $\max$ preserves KP summability with the constants of~\eqref{eq:buffer-kp-constants}) is settled in the present document; what 3.3 owes is the regime-specific verification that the resulting $(\rho_{\cap}, K_{\cap})$ is tight enough to feed the inductive step of Milestone~3.4.
\end{itemize}
The role of the present document is to fix the predicate's signature and to establish the structural Proposition~\ref{prop:summability-preservation}; the role of Milestones 3.1--3.3 is to exhibit concrete witnesses.
\end{remark}

\section{Domination gluing on the buffer overlap}
\label{sec:gluing}

\begin{definition}[Domination gluing]
\label{def:domGlue}
Given an admissible small-field bound $s^{\mathrm{sf}} \in \Acal_k(\Omsf, X)$ and an admissible large-field bound $s^{\mathrm{lf}} \in \Acal_k(\Omlf, X)$, the \emph{domination glue} on the buffer is
\begin{equation}
  \label{eq:domGlue}
  s^{\cap}(u) \;:=\; \max\!\big( s^{\mathrm{sf}}(u),\ s^{\mathrm{lf}}(u) \big)
  \qquad \text{for } u \in \Omcap.
\end{equation}
(Where one of $s^{\mathrm{sf}}, s^{\mathrm{lf}}$ is not defined at $u$, we take $\max$ to be the other; in the presentation chosen in this document both are everywhere-defined on $\Omk$ but only admissible on their respective region.)
\end{definition}

\begin{proposition}[Domination gluing is admissible]
\label{prop:domGlue-admissible}
Under the standing assumption of Definition~\ref{def:activity}, the domination glue~\eqref{eq:domGlue} is in $\Acal_k(\Omcap, X)$.
\end{proposition}

\begin{proof}
For $u \in \Omcap$, the activity $|a_k(X)|(u)$ is a single nonneg real (Definition~\ref{def:activity}).  The pointwise max satisfies $s^{\cap}(u) \ge s^{\mathrm{sf}}(u)$ and $s^{\cap}(u) \ge s^{\mathrm{lf}}(u)$ trivially.  Now,~$\Omcap \subseteq \Omsf \cup \Omlf$ (every buffer configuration is in either the small-field or the large-field region, possibly both), so for every $u \in \Omcap$ at least one of $|a_k(X)|(u) \le s^{\mathrm{sf}}(u)$ or $|a_k(X)|(u) \le s^{\mathrm{lf}}(u)$ holds; hence $|a_k(X)|(u) \le s^{\cap}(u)$.
\end{proof}

\begin{remark}[Why domination, not equality]
Two distinct upper bounds $s^{\mathrm{sf}}$ and $s^{\mathrm{lf}}$ on the same activity will not in general agree on the overlap: the small-field bound is derived from a Gaussian/perturbative expansion of the action, the large-field bound from a heat-kernel character expansion (outline \S6.4, \S6.5).  They bound the same function from above, but disagree pointwise on $\Omcap$.  Gluing by equality would fail.  Gluing by domination (Proposition~\ref{prop:domGlue-admissible}) absorbs the disagreement by taking the looser bound on the overlap.  This is the structural pay-off of the sheaf reformulation announced in \texttt{docs/context.md} \S3.1: the cocycle problem becomes a domination problem.
\end{remark}

\section{Summability preservation under domination gluing}
\label{sec:buffer-summability}

Proposition~\ref{prop:domGlue-admissible} establishes that the
domination glue is an admissible \emph{bound}.  The trivial section
$\infty$ is also admissible (Theorem~\ref{thm:flabbiness}), so
admissibility alone has no analytic content.  The non-trivial half
of the buffer mechanism is the question: \emph{does the domination
glue preserve KP summability, and with what constants?}  This
section states the load-bearing claim that answers it.

\begin{proposition}[Summability preservation under domination gluing]
\label{prop:summability-preservation}
Let $X \in \Pcal_k$ be a polymer.  Suppose:
\begin{enumerate}
  \item[(i)] A coherent family $\{s^{\mathrm{sf}}_{X'}\}_{X' \supseteq X}$ of bound-witness sections on $\Omsf$ is compatible with the KP datum $(\rho_{\mathrm{sf}}, K_{\mathrm{sf}}) \in \Scal_k(\Omsf, X)$ in the sense of Definition~\ref{def:compatibility}.
  \item[(ii)] A coherent family $\{s^{\mathrm{lf}}_{X'}\}_{X' \supseteq X}$ of bound-witness sections on $\Omlf$ is compatible with $(\rho_{\mathrm{lf}}, K_{\mathrm{lf}}) \in \Scal_k(\Omlf, X)$.
\end{enumerate}
Define the buffer-glued family $\{s^{\cap}_{X'}\}_{X' \supseteq X}$ on $\Omcap$ by Definition~\ref{def:domGlue} (pointwise maximum, falling back to the single available bound where only one of $s^{\mathrm{sf}}, s^{\mathrm{lf}}$ is defined).  Then the buffer family is compatible with the KP datum
\begin{equation}
  \label{eq:buffer-kp-constants}
  (\rho_{\cap}, K_{\cap}) \;:=\; \bigl(\, \min(\rho_{\mathrm{sf}}, \rho_{\mathrm{lf}}),\ \ K_{\mathrm{sf}} + K_{\mathrm{lf}}\, \bigr) \;\in\; \Scal_k(\Omcap, X).
\end{equation}
\end{proposition}

\begin{proof}[Proof sketch]
Fix $u \in \Omcap$ and a plaquette $x$ of $X$.  By the cover structure (Definition~\ref{def:three-cover}, $\Omcap \subseteq \Omsf \cup \Omlf$), $u$ lies in at least one of the two regions; we split into two cases.

\textbf{Case A:} $u \in \Omsf \cap \Omlf$.  Both $s^{\mathrm{sf}}_{X'}(u)$ and $s^{\mathrm{lf}}_{X'}(u)$ are defined for every $X' \supseteq X$, and $s^{\cap}_{X'}(u) = \max\!\bigl(s^{\mathrm{sf}}_{X'}(u), s^{\mathrm{lf}}_{X'}(u)\bigr)$.  Applying $\max(a, b) \le a + b$:
\begin{align*}
  \sum_{X' \supseteq X,\ X' \ni x} s^{\cap}_{X'}(u)\, e^{\rho_{\cap} |X'|}
  &\le\; \sum_{X'} s^{\mathrm{sf}}_{X'}(u)\, e^{\rho_{\cap} |X'|} \;+\; \sum_{X'} s^{\mathrm{lf}}_{X'}(u)\, e^{\rho_{\cap} |X'|}.
\end{align*}
Since $\rho_{\cap} \le \rho_{\mathrm{sf}}$ by~\eqref{eq:buffer-kp-constants}, $e^{\rho_{\cap} |X'|} \le e^{\rho_{\mathrm{sf}} |X'|}$, and by hypothesis~(i) the first sum is $\le K_{\mathrm{sf}}$.  Symmetrically the second is $\le K_{\mathrm{lf}}$.  Total: $\le K_{\mathrm{sf}} + K_{\mathrm{lf}} = K_{\cap}$.

\textbf{Case B:} $u \in \Omlf \setminus \Omsf$ (or symmetrically $u \in \Omsf \setminus \Omlf$).  Only one of $s^{\mathrm{sf}}, s^{\mathrm{lf}}$ is defined at $u$; by Definition~\ref{def:domGlue} $s^{\cap}_{X'}(u)$ equals that single value.  The KP sum is then $\le K_{\mathrm{lf}} \le K_{\mathrm{sf}} + K_{\mathrm{lf}}$ (resp.\ $\le K_{\mathrm{sf}}$).
\end{proof}

\begin{remark}[Looseness of the union-bound estimate]
\label{rem:union-bound-looseness}
The step $\max(a, b) \le a + b$ in Case~A is the union bound, loose by at most a factor of $2$.  \texttt{docs/context.md} \S3.4 raises a parallel concern about a \emph{different} union bound (the per-plaquette union bound on the overlap \emph{measure}, not the KP-summability glue treated here).  For the present statement, the looseness factor is constant in $k$, so it does not accumulate across RG steps; the substantive tightness question is whether $\rho_{\cap} = \min(\rho_{\mathrm{sf}}, \rho_{\mathrm{lf}})$ stays bounded away from zero across the RG trajectory.  That is the content of Milestone~3.4 (string-tension 2-cell budget), not of this proposition.
\end{remark}

\begin{remark}[Discharge schedule]
\label{rem:summability-discharge}
Proposition~\ref{prop:summability-preservation} is the load-bearing claim of the buffer mechanism.  Its proof above is structural and assumes only the standing setup of the framework (Definition~\ref{def:three-cover} cover structure, Definition~\ref{def:compatibility} compatibility).  The proposition is therefore stated and proved \emph{here}, at the framework level (Milestone~2.1), rather than deferred entirely to Milestone~3.3 as the previous version of this document did (see \texttt{docs/sheaf-framework.md} \S5).  Milestone~3.3's remaining obligation is the \emph{regime-specific} verification: producing concrete coherent families $\{s^{\mathrm{sf}}_{X'}\}$ and $\{s^{\mathrm{lf}}_{X'}\}$ for $SU(2)$ at each scale~$k$, with explicit $(\rho_{\mathrm{sf}}, K_{\mathrm{sf}})$ and $(\rho_{\mathrm{lf}}, K_{\mathrm{lf}})$, and showing the resulting $(\rho_{\cap}, K_{\cap})$ from~\eqref{eq:buffer-kp-constants} is small enough to feed the inductive step (Milestone~3.4).
\end{remark}

\begin{conjecture}[Expected sharpness of the union-bound estimate, to be verified]
\label{conj:gaussian-sharpness}
The author's judgment, stated as a \emph{conjecture} and not a proved fact, to be tested against the SU(2) constants computed in Milestones 1.3, 1.4, and 3.1--3.3:
\begin{quote}
For heat-kernel / Gaussian polymer activities, the union-bound estimate $K_{\cap} \le K_{\mathrm{sf}} + K_{\mathrm{lf}}$ is sharp up to a constant factor $\le 2$.  The expected mechanism: for each polymer $X'$, exactly one of $s^{\mathrm{sf}}_{X'}$ or $s^{\mathrm{lf}}_{X'}$ saturates the maximum --- the small-field bound dominates for short polymers, the large-field bound for long ones --- so the polymer sum splits naturally into ``mostly sf'' and ``mostly lf'' contributions and the union bound is rarely a true factor-of-$2$ lossy.
\end{quote}

\textbf{Unit-mismatch caveat.}  The above expectation is \emph{not} a settled fact because the two regime bounds live in different units: the small-field activity is bounded by $\exp(-\rho_{\mathrm{sf}} |X'|_{\mathrm{phys}}/a_k)$ (continuum-like, indexed by physical extent in lattice units) while the large-field activity is bounded by $\exp(-\rho_{\mathrm{lf}} |X'|_{\mathrm{plaq}})$ (lattice-combinatorial, indexed by plaquette count).  Whether these two regimes admit a clean crossover at a polymer-by-polymer level --- so that the conjecture's ``exactly one saturates'' claim holds --- or whether there is a fat buffer band of polymers for which neither bound is tight is \emph{exactly} the unit-mismatch difficulty that \texttt{docs/context.md} \S3.1 names as the original obstruction the sheaf reformulation is built to absorb.  Until concrete forms of $s^{\mathrm{sf}}_{X'}$, $s^{\mathrm{lf}}_{X'}$ exist (Milestones 1.3, 1.4) and the resulting union bound is checked against direct estimation on the buffer (Milestone 3.3), the conjecture stands as an expectation, not a result.  This is the same flavour of warning \texttt{docs/context.md} \S3.4 raises about the parallel per-plaquette union bound for the overlap \emph{measure}: union bounds in this framework ``do real work and may be too loose,'' and tightness has to be checked against concrete polymer-activity forms, not assumed.

A softer glue does not help here.  Replacing $\max$ with log-sum-exp on linear values gives $\log(s^{\mathrm{sf}} + s^{\mathrm{lf}})$ in linear scale, which is the union bound itself.  A convex combination $\lambda s^{\mathrm{sf}} + (1-\lambda) s^{\mathrm{lf}}$ with $\lambda \in (0,1)$ is smaller than $\max$ and would fail admissibility (Proposition~\ref{prop:domGlue-admissible}), undoing the structural pay-off of the sheaf reformulation.  The right glue for this framework is $\max$, and the buffer mechanism's tightness is controlled by the $\rho$-degradation budget across RG steps and by whether the conjecture above turns out to hold, not by the form of the glue itself.

If the inductive step fails to close at Milestone~3.4, the author's expectation is that the failure will surface either in the $\rho_{\cap}$ trajectory ($\min(\rho_{\mathrm{sf}}, \rho_{\mathrm{lf}})$ degrading too fast as $k$ grows) or in the conjecture above turning out to be false on a fat buffer band of polymers.  The structural ($K_{\cap}$) framework is robust; both robustness questions are analytic and owed to Phase~3.
\end{conjecture}

\section{Lean mirror}
\label{sec:lean}

The Lean formalization of this document lives in \texttt{lean/YMGap/Sheaf/}:
\begin{itemize}
  \item \texttt{Site.lean} --- the configuration site (\texttt{Open}, \texttt{Object}, \texttt{Morphism}, \texttt{ThreeCover}; cf.\ Section~\ref{sec:site}).
  \item \texttt{BoundWitness.lean} --- the bound-witness presheaf abstract scaffold (\texttt{BoundValues} typeclass, \texttt{Section}, the flabbiness theorem; cf.\ Section~\ref{sec:awitness}).
  \item \texttt{Summability.lean} --- the summability-certificate record (\texttt{KPCertificate}; cf.\ Section~\ref{sec:scertif}).
  \item \texttt{Gluing.lean} --- the domination glue (\texttt{domGlue}) and the named compatibility predicate (\texttt{CompatibilityRule}; cf.\ Sections~\ref{sec:compat} and~\ref{sec:gluing}).
\end{itemize}
The value type for $\Acal_k$-sections in Lean is intentionally abstract: a generic type $T$ with $[\mathrm{BoundValues}\ T]$, where $\mathrm{BoundValues}$ records the partial order, the top element $\top$, and the binary sup.  The math paper instantiates $T = \Rnnext$; the Lean formalization defers the value-type commitment until \texttt{lean/YMGap/Constructive/Reals.lean} is built (currently not implemented).  Per AGENTS.md Rule~3 and the empirical finding of Milestone~1.1 that Lean~4 core's \texttt{Rat} drags \texttt{Classical.choice} into reflexive equality, the Lean value type cannot be a real-number type built on Mathlib without breaking the axiom-clean discipline.

The structural lemmas of this document --- Theorem~\ref{thm:flabbiness} (flabbiness) and Proposition~\ref{prop:domGlue-admissible} (domination gluing) --- are mirrored in Lean at the abstract value-type level, with empty axiom dependency sets.

\section{Cross-references}
\label{sec:xref}

\begin{itemize}
  \item Prose companion: \texttt{docs/sheaf-framework.md}.
  \item Lean formalization: \texttt{lean/YMGap/Sheaf/}.
  \item Diagram: \texttt{visualizations/structure/sheaf-gluing.tikz} (the three-region cover with the coupled $\Acal_k$, $\Scal_k$ sheaves).
  \item Motivation: \texttt{docs/context.md} \S3.1 (cocycle observation), \S3.2 (sheaf reformulation), \S3.4 (overlap measure bound).
  \item Specification: \texttt{docs/outline.md} \S8 (sheaf framework axiomatics) and \S9 (buffer-overlap estimate).
  \item Downstream consumers: Milestones 2.2 (inductive state), 2.3 (block-spin morphism), 3.1--3.3 (regime KP constants), 3.4 (string-tension 2-cell).
\end{itemize}

\end{document}
